% !TeX root = main.tex
\section{Method: RSVP}

\subsection{Predicate-Wise Residual Shaping Value}

RSVP는 기본 MARL 학습기의 가치 함수와 별도로 보조 예측기 하나를 둔다. 먼저
$x_t=\psi(s_t)$를 갱신 판단에 필요한 과업 상태의 저차원 수치 요약이라 하자.
$x_t$는 LLM 입력용 텍스트 $d_t$와 구별되며, 현재 지침의 문장이나 가중치를 포함하지
않는다. 이 수치 요약을 입력받아 $M$개 predicate의 값을 함께 예측하는 회귀 모델을
이하 \emph{monitoring critic}이라 부른다. 여기서 critic은 저수준 행동을 선택하는
action-value critic이 아니라 갱신 판단만을 위한 예측기이다. 각 출력 성분을 predicate
$j$의 head라 부른다.

head $j$가 예측하는 residual shaping value는 시점 $t$부터 episode가 끝날 때까지
predicate $j$가 얼마나 더 발생할지를 나타내는 할인 누적량이다. 완료된 episode에서
종료 직후의 인덱스를 $T$라 하면 학습 타깃은

\begin{equation}
Y_{j,t}=\sum_{\ell=t}^{T-1}\gamma_F^{\ell-t}f_{j,\ell}
=f_{j,t}+\gamma_FY_{j,t+1},\qquad Y_{j,T}=0
\label{eq:critic-target}
\end{equation}

이다. $\gamma_F\in[0,1)$는 모니터가 바라보는 시간 범위를 정하며, 환경 보상의
할인율과 별도로 설정할 수 있다. monitoring critic의 $j$번째 출력
$\widehat V_j(x_t)$은 $Y_{j,t}$를 회귀한다. 활성 집합 $\mathcal I_r$ 가운데 타깃
분산이 충분해 학습 가능하다고 판단된 head의 집합을 $\mathcal H_r$라 하면, 현재
지침과 관련된 예측값은

\begin{equation}
\widehat V_F(x;G_r)=
\sum_{j\in\mathcal H_r}w_{r,j}\widehat V_j(x)
\label{eq:guidance-value}
\end{equation}

로 합성한다. 이 값은 식~\eqref{eq:shaping}의 clip이나 식~\eqref{eq:reward}의
$\lambda_t$를 적용하기 전의 가중 합이다. predicate별로 출력을 분리했기 때문에 같은
monitoring critic을 서로 다른 지침에 재사용할 수 있다. 이 논문에서 \emph{residual}은
현재 시점 뒤에 남은 누적량이라는 뜻이며 TD residual을 뜻하지 않는다. 또한
$\widehat V_F$는 외부 return이나 지침을 유지했을 때의 반사실적 이득이 아니다. 현재
정책 아래에서 지침에 대응하는 측정 신호가 더 발생할 가능성을 근사한다.

\subsection{Issuance Guard and Cumulative Degradation Detection}

새 지침이 발급된 시점 $r$의 수치 특징을 $x_r$라 하자. 온라인 초기의 임의 예측이나
거의 항상 0인 predicate로 갱신이 발생하는 것을 막기 위해 발급 안전장치(issuance
guard)를 둔다. 구체적으로, 모니터 학습 버퍼에서 $Y_j$의 경험적 분산이 최소값보다
큰 head만 학습 가능하다고 간주한다. 활성 가중치 중 이러한 head에 연결된 비율

\begin{equation}
q_r=\frac{\sum_{j\in\mathcal H_r}w_{r,j}}
{\sum_{(j,w_{r,j})\in\mathcal I_r}w_{r,j}}
\label{eq:coverage}
\end{equation}

을 \emph{coverage}라 한다. 또한 버퍼에서 타깃 $Y_j$의 경험적 평균을 $\mu_j$라 하고,
지침의 기준 크기를

\begin{equation}
b_r=\sum_{j\in\mathcal H_r}w_{r,j}\max(\mu_j,0)
\label{eq:issuance-base}
\end{equation}

로 계산한다. 수집한 episode 수가 $n_{\min}$ 이상이고, $q_r\ge q_{\min}$이며,
$b_r\ge b_{\min}$이고, 발급 시 예측값이
$\widehat V_F(x_r;G_r)\ge\epsilon_I b_r$일 때만 조기 검출기를 활성화한다. 이 조건을
통과하지 못하면 해당 지침에는 최대 갱신 간격만 적용한다.

활성 구간에서는 현재 critic 파라미터로 현재 상태와 발급 상태를 함께 재평가하여

\begin{equation}
v_t=\min\!\left(1,
\frac{\max(0,\widehat V_F(x_t;G_r))}
{\max(\widehat V_F(x_r;G_r),\epsilon)}
\right)
\label{eq:ratio}
\end{equation}

를 계산한다. 분모도 발급 당시의 오래된 파라미터로 고정하지 않고, 매 스텝 현재
monitoring critic으로 저장된 $x_r$를 다시 평가한다. 이는 학습 중 모델 파라미터가
변한 효과가 비율에 직접 섞이는 것을 줄인다. $\epsilon>0$은 0에 가까운 분모를 막는
상수이다. 정상적으로 남아 있다고 간주하는 비율을 $\kappa$, 검출 임계값을 $h$라 하고
단측 누적합 통계량~\cite{page1954cusum}

\begin{equation}
S_t=\max(0,S_{t-1}+\kappa-v_t)
\label{eq:cusum}
\end{equation}

을 갱신한다. $v_t<\kappa$인 상태가 이어지면 $S_t$가 증가한다. 최소 보유 간격을
지난 뒤 $S_t\ge h$가 되면 Commander를 조기
호출한다. 호출에 성공하면 새 지침과 기준 상태를 저장하고 $S_t$를 초기화한다.
마지막 성공 호출 후 $F_{\max}$스텝이 지나면 검출 여부와 관계없이 강제 호출을
시도한다. 따라서 $F_{\max}$는 정상 호출 조건에서 지침 보유 시간의 상한이며,
$h$와 $\kappa$는 조기 호출의 민감도를 조절한다.

CUSUM은 단일 스텝의 희소하고 변동성 큰 predicate 값에 반응하는 대신, 예측된 감소가
지속될 때 증거를 누적하기 위해 사용한다. 그러나 식~\eqref{eq:cusum}은 알려진 두
확률분포의 로그 우도비로 유도된 것이 아니다. 비정상적인 학습 과정에서 사용하는
경험적 변화 감지 규칙이므로 고전적 CUSUM의 최적성을 주장하지 않는다.

\subsection{Training and Computational Flow}

각 환경 스텝에서 RSVP는 (1) $x_t$를 계산하고, (2) 조기 또는 강제 갱신 여부를
판단하고, (3) 모든 접지 함수의 출력과 $F_t$를 계산하며, (4) 전이를 기본 MARL
학습기의 재현 버퍼에 저장한다. Episode가 끝나면 식~\eqref{eq:critic-target}의 타깃을
별도의 모니터 학습 버퍼에 추가하고 monitoring critic을 학습한다. Commander 호출
사이에는 기존 지침을 재사용한다. 감시에는 접지 과정에서 이미 계산 가능한 신호를
이용하므로 별도의 LLM 판정 호출이 필요하지 않지만, monitoring critic의 추론과
학습에 따른 추가 계산은 발생한다.
